\documentclass[11pt,a4paper]{article}
\usepackage[utf8]{inputenc}
\usepackage[T1]{fontenc}
\usepackage{geometry}
\geometry{margin=1in}
\usepackage{hyperref}
\usepackage{enumitem}
\usepackage{booktabs}
\usepackage{longtable}
\usepackage{xcolor}
\usepackage{titlesec}
\usepackage{amsmath,amssymb}

\titleformat{\section}{\Large\bfseries\color{blue!70!black}}{}{0em}{}
\titleformat{\subsection}{\large\bfseries\color{blue!50!black}}{}{0em}{}
\titleformat{\subsubsection}{\normalsize\bfseries\color{blue!30!black}}{}{0em}{}

\title{\textbf{Replication Plan}\\[0.5em]
\large Physics and Chemistry from Parsimonious Representations: Image Analysis via Invariant Variational Autoencoders\\[0.3em]
\normalsize OSTI ID: 2439897}
\author{Auto-generated Replication Blueprint}
\date{\today}

\begin{document}
\maketitle
\tableofcontents
\newpage

%---------------------------------------------------------------------
\section{Paper Overview}
%---------------------------------------------------------------------
This paper presents invariant variational autoencoders (iVAEs) for learning disentangled, physically meaningful latent representations from microscopy and spectroscopy images. The approach builds translational and rotational invariance into the VAE architecture, enabling the model to separate physical degrees of freedom (e.g., composition, strain, defect concentration) from nuisance variables (position, orientation). The authors apply the method to scanning probe microscopy (SPM) and electron microscopy data from materials science, demonstrating that the latent space organizes according to physical parameters.

%---------------------------------------------------------------------
\section{Detailed Workflow}
%---------------------------------------------------------------------

\subsection{Phase 1: Data Preparation}
\begin{enumerate}[label=\textbf{1.\arabic*}]
  \item \textbf{Obtain microscopy image datasets.}
    \begin{itemize}
      \item Simulated STEM or SPM images with known physical parameters (for validation).
      \item Experimental microscopy images from the authors or public repositories.
    \end{itemize}
  \item \textbf{Preprocess images.}
    \begin{itemize}
      \item Normalize intensities, crop/resize to uniform dimensions.
      \item Extract patches (sub-images) centered on atomic positions or grid-sampled.
      \item Create train/validation/test splits.
    \end{itemize}
\end{enumerate}

\subsection{Phase 2: iVAE Architecture}
\begin{enumerate}[label=\textbf{2.\arabic*}]
  \item \textbf{Implement the encoder.}
    \begin{itemize}
      \item Convolutional layers extracting features from image patches.
      \item Rotational invariance: use circular harmonics, steerable CNNs, or data augmentation.
      \item Translational invariance: global average pooling or spatial transformer.
      \item Output: mean $\mu$ and log-variance $\log \sigma^2$ of the latent distribution.
    \end{itemize}
  \item \textbf{Implement the decoder.}
    \begin{itemize}
      \item Transposed convolutions reconstructing the image from the latent code.
    \end{itemize}
  \item \textbf{Loss function.}
    \begin{itemize}
      \item Reconstruction loss (MSE or BCE) + KL divergence regularization.
      \item Optional: $\beta$-VAE weighting for disentanglement.
    \end{itemize}
\end{enumerate}

\subsection{Phase 3: Training}
\begin{enumerate}[label=\textbf{3.\arabic*}]
  \item Train the iVAE on the image patch dataset.
  \item Hyperparameter tuning: latent dimension, $\beta$, learning rate, architecture depth.
  \item Monitor reconstruction quality and KL divergence.
\end{enumerate}

\subsection{Phase 4: Latent Space Analysis}
\begin{enumerate}[label=\textbf{4.\arabic*}]
  \item \textbf{Visualize the latent space} (2D projections via t-SNE or UMAP).
  \item \textbf{Correlate latent dimensions} with known physical parameters (composition, strain).
  \item \textbf{Latent space traversals}: vary one latent dimension and decode to visualize the learned transformation.
  \item \textbf{Clustering in latent space} to identify distinct phases or structures.
\end{enumerate}

\subsection{Phase 5: Validation and Figure Reproduction}
\begin{enumerate}[label=\textbf{5.\arabic*}]
  \item Reproduce latent space organization plots.
  \item Reproduce reconstruction examples.
  \item Reproduce quantitative metrics (disentanglement scores, reconstruction error).
\end{enumerate}

%---------------------------------------------------------------------
\section{Datasets}
%---------------------------------------------------------------------

\begin{longtable}{p{4.5cm} p{3.5cm} p{6cm}}
\toprule
\textbf{Dataset / Source} & \textbf{Type} & \textbf{Location / Notes} \\
\midrule
\endfirsthead
\toprule
\textbf{Dataset / Source} & \textbf{Type} & \textbf{Location / Notes} \\
\midrule
\endhead
Simulated microscopy images & Synthetic & Generated via multislice or analytical models \\
\midrule
Experimental SPM/STEM data & From authors & Check paper supplementary or ORNL data repositories \\
\midrule
Code repository data & Bundled & \url{https://github.com/} (check paper for specific repo) \\
\midrule
dSprites (baseline benchmark) & Public & \url{https://github.com/deepmind/dsprites-dataset} \\
\bottomrule
\end{longtable}

%---------------------------------------------------------------------
\section{Models, Tools, and Software}
%---------------------------------------------------------------------

\begin{longtable}{p{4.5cm} p{2.5cm} p{6.5cm}}
\toprule
\textbf{Tool / Library} & \textbf{Version} & \textbf{Purpose / URL} \\
\midrule
\endfirsthead
\toprule
\textbf{Tool / Library} & \textbf{Version} & \textbf{Purpose / URL} \\
\midrule
\endhead
PyTorch & $\geq 1.10$ & Deep learning framework \\
\midrule
NumPy / SciPy & $\geq 1.21$ & Numerical operations \\
\midrule
Matplotlib & $\geq 3.4$ & Visualization \\
\midrule
scikit-learn & $\geq 1.0$ & Clustering, dimensionality reduction \\
\midrule
UMAP-learn & latest & UMAP visualization \newline \url{https://umap-learn.readthedocs.io/} \\
\midrule
AtomAI & latest & Deep learning for microscopy (ORNL) \newline \url{https://github.com/pycroscopy/atomai} \\
\midrule
e2cnn & latest & Equivariant CNNs for rotational invariance \newline \url{https://github.com/QUVA-Lab/e2cnn} \\
\midrule
Jupyter & latest & Notebook environment \\
\bottomrule
\end{longtable}

\subsection{Hardware Requirements}
\begin{itemize}
  \item GPU with $\geq$ 8\,GB VRAM for training.
  \item $\geq$ 16\,GB RAM.
\end{itemize}

\end{document}
